\documentclass[aspectratio=169,11pt]{beamer}

\usetheme{Madrid}
\usecolortheme{seahorse}
\setbeamertemplate{navigation symbols}{}
\setbeamertemplate{footline}[frame number]

\usepackage[T1]{fontenc}
\usepackage{booktabs}
\usepackage{array}
\usepackage{xcolor}
\usepackage{listings}
\usepackage{amsmath}
\usepackage{amssymb}
\usepackage{hyperref}

\definecolor{codebg}{gray}{0.95}
\lstset{
  basicstyle=\ttfamily\scriptsize,
  backgroundcolor=\color{codebg},
  breaklines=true,
  columns=fullflexible,
  showstringspaces=false,
  frame=single
}

\title{NixNan: Detecting Floating-Point Exceptions on GPUs}
\subtitle{Architecture, Features, and Four Case Studies}
\author{SC26 Tutorial}
\date{}

\begin{document}

\begin{frame}
\titlepage
\end{frame}

% ---------------------------------------------------------------
\section{Motivation}

\begin{frame}[fragile]{Why Floating-Point Exceptions Matter}
\begin{itemize}
  \item FP exceptions (NaN, Inf) are \textbf{almost always bugs}, not intended results.
  \item They are easy to generate, especially in dot products / reductions, and
    especially in \textbf{low-precision ML formats}:
  \[
  (1.9\times10^{19})^2 \to \mathtt{INF} \text{ in FP32}, \qquad
  16^2 = \mathtt{INF} \text{ in E4M3 (a popular ML format)}
  \]
  \item A classic, deceptively simple bug:
\begin{lstlisting}[language=C]
#define MAX(x,y)  ( (x >= y) ? x : y )
\end{lstlisting}
  With $x=\mathtt{NaN}, y=42$: result is $42$ (looks clean). Swap the
  arguments: result is $\mathtt{NaN}$. \textbf{MAX must be commutative --- it isn't.}
  \item NaNs can silently skew \textbf{control flow}, not just numeric output.
\end{itemize}
\end{frame}

\begin{frame}{The Silent-Killer Problem, and Why GPUs Make It Worse}
\begin{itemize}
  \item Exceptions often \textbf{vanish harmlessly} on the way to program output
    (masked by a later operation, a reduction, a clamp) --- but it is very
    hard to know, after the fact, \textbf{which control-flow decisions they
    influenced} while they were alive.
  \item Results are only easy to observe at the very \emph{end} of a kernel
    or program. If an exception dies before reaching there, the programmer
    never sees it.
  \item NVIDIA GPUs generate \textbf{no hardware traps} on FP exceptions ---
    unlike a CPU's SIGFPE.
  \item So exception detection must happen \textbf{in software}, by decoding
    the bit pattern of every computed value:
  \[
  \underbrace{S}_{\text{sign}}\;\underbrace{E\cdots E}_{\text{exponent}}\;\underbrace{F\cdots F}_{\text{mantissa}}
  \]
  \item All-ones exponent + nonzero mantissa $=$ NaN; all-ones exponent + zero
    mantissa $=$ $\pm\infty$; all-zeros exponent + nonzero mantissa $=$ subnormal.
  \item This has to happen for \textbf{every FP instruction, every thread,
    every launch} --- at native GPU execution speed, or close to it.
\end{itemize}
\end{frame}

% ---------------------------------------------------------------
\section{NixNan: Architecture}

\begin{frame}{From GPU-FPX to NixNan}
\begin{columns}
\column{0.5\textwidth}
\textbf{GPU-FPX} (HPDC 2023)
\begin{itemize}
  \item Separate detector + analyzer (flow analysis)
  \item Exits early on first exception per thread
  \item No FP16 or Tensor Core support
\end{itemize}
\column{0.5\textwidth}
\textbf{NixNan} (evolved from GPU-FPX)
\begin{itemize}
  \item Single component: detects \emph{and} flow-analyzes
  \item Finds \emph{all} exceptions --- never quits early
  \item Full FP16 and Tensor Core support
\end{itemize}
\end{columns}
\vspace{0.5em}
\begin{itemize}
  \item Continuously refined against real workloads --- this tutorial's own
    four case studies found and fixed two more bugs in NixNan itself
    (Slide~\ref{sl:dogfood}).
\end{itemize}
\end{frame}

\begin{frame}{Where NixNan Hooks Into the Compilation Pipeline}
\begin{center}
\begin{tabular}{ccccc}
CUDA C code & $\xrightarrow{\texttt{nvcc}}$ & PTX code & $\xrightarrow{\texttt{ptxas}}$ & SASS code \\
\end{tabular}
\end{center}
\begin{itemize}
  \item Small matrix multiplications $D = A{\cdot}B + C$ are performed
    \textbf{directly by Tensor Core hardware} --- not expressed as ordinary
    scalar FP instructions at all.
  \item NixNan instruments \textbf{SASS} (the actual machine code), not the
    source or PTX: it sees exactly what runs on the GPU, \emph{after} every
    compiler optimization.
  \item Consequence: \textbf{closed/precompiled libraries can be
    ``X-rayed''} too --- cuBLAS, cuSOLVER kernels included, no source needed.
\end{itemize}
\end{frame}

\begin{frame}[fragile]{How NixNan Instruments: NVBit + \texttt{LD\_PRELOAD}}
\begin{itemize}
  \item NixNan is an \textbf{NVBit} (NVIDIA Binary Instrumentation Tool) tool.
\begin{lstlisting}[language=bash]
LD_PRELOAD=nixnan.so  ./your_precompiled_cuda_app
\end{lstlisting}
  \item NVBit intercepts every CUDA driver \texttt{Kernel Launch} call and
    gives NixNan a callback \emph{before} that kernel's SASS runs.
  \item The SASS rewrite (inserting the checking calls) happens only the
    \textbf{first time} a given kernel function is seen --- later launches of
    the same kernel reuse the already-instrumented code, so the rewrite cost
    is paid once, not per launch.
  \item Detects exceptions in \textbf{both SIMT cores and Tensor Cores}.
\end{itemize}
\end{frame}

\begin{frame}[fragile]{Detecting Exceptions: Bit-Pattern Classification}
\begin{itemize}
  \item For every FP-classified instruction, NixNan inserts a call to a
    device-side checker \emph{before} and \emph{after} it executes.
  \item Classification is pure bit-masking on the operand's exponent/mantissa
    --- \textbf{no floating-point arithmetic}, so the check itself cannot
    introduce rounding artifacts:
\begin{lstlisting}[language=C]
uint32_t float_is_nan(uint32_t f) {
  return (float_exp(f)==0xFF && float_mant(f)!=0) ? E_NAN : 0;
}
\end{lstlisting}
  \item Per-lane exception bits are reduced across the warp
    (\texttt{\_\_shfl\_sync}); only one lane per warp records, avoiding
    32-way redundant writes.
  \item First occurrence of a given (instruction, exception-type) pair is
    pushed live through a device$\to$host channel to a receiver thread.
\end{itemize}
\end{frame}

\begin{frame}{What NixNan Detects}
\begin{center}
\begin{tabular}{ll}
\toprule
Category & Meaning \\
\midrule
NaN & Not-a-number \\
Infinity / $-$Infinity & Overflow (reported with sign) \\
Subnormal & Denormalized (near-zero) value \\
Division by zero & Divisor operand exactly zero \\
\bottomrule
\end{tabular}
\end{center}
\begin{itemize}
  \item Across \textbf{FP16, BF16, FP32, FP64} (FP8 is future work).
  \item Across \textbf{both SIMT-core scalar instructions} (FADD, FFMA, DFMA, \ldots)
    \textbf{and Tensor Core MMA instructions} (HMMA, confirmed empirically ---
    see Case Study~3).
  \item Optional \texttt{INSTR\_MEM=1}: also instrument global-memory
    \textbf{stores}, to confirm an exception was actually \emph{written},
    not just transiently computed.
\end{itemize}
\end{frame}

\begin{frame}{Key Environment Variables}
\scriptsize
\begin{center}
\begin{tabular}{lp{8cm}}
\toprule
Variable & Effect \\
\midrule
\texttt{LD\_PRELOAD} & Points to \texttt{nixnan.so}; how the tool attaches at all \\
\texttt{INSTRUMENT\_EXCEPTIONS} & Master switch for NaN/Inf/Subnormal/Div0 checking (default \textbf{on}) \\
\texttt{INSTR\_MEM} & Also check global-memory stores \\
\texttt{HISTOGRAM} / \texttt{BIN\_SPEC\_FILE} & Exponent-range (binade) histogramming, Slide~\ref{sl:histogram} \\
\texttt{SAMPLING} & Instrument only every $N$th launch of a repeated kernel \\
\texttt{MAX\_ERRORS} & Self-terminate after this many exception reports \\
\texttt{LOG\_KERNELS} & \textbf{Disables} all instrumentation; just logs kernel launches + timing \\
\texttt{TOOL\_VERBOSE}, \texttt{ENABLE\_FUN\_DETAIL} & Extra diagnostic printing at setup / per-launch \\
\texttt{FUNCTION\_WHITELIST} / \texttt{\_BLACKLIST} & Restrict instrumentation to (or exclude) named kernels \\
\bottomrule
\end{tabular}
\end{center}
\end{frame}

\begin{frame}{The Exponent-Histogram Mechanism}
\label{sl:histogram}
\begin{itemize}
  \item \texttt{BIN\_SPEC\_FILE}: a JSON spec naming, per format, exponent
    ranges to watch, a report \texttt{count}, and (optionally) \texttt{doublings}.
  \item Every FP value updates a \texttt{[format][exponent]} histogram
    (powers the \texttt{"Exponent range for f16: [-9, 8]"} summary line).
  \item If the exponent falls in a watched bin, a per-bin counter increments;
    every \texttt{count}-th hit sends a \textbf{threshold report}.
  \item \textbf{A real blind spot}: valid bins explicitly exclude the zero and
    infinity exponent codes. \textbf{A value that jumps straight from finite
    to Inf in one instruction never lands in any bin} --- invisible to this
    mechanism (see Slide~\ref{sl:blindspot}).
\end{itemize}
\end{frame}

\begin{frame}{Adaptive Doubling: an Exponentially Decaying Report Rate}
\begin{itemize}
  \item Reporting on \emph{every} occurrence at \texttt{count:1} floods the
    log for a hot kernel. \texttt{doublings} fixes this:
  \begin{itemize}
    \item Start reporting every \texttt{count} occurrences.
    \item Each report \textbf{doubles} that bin's own threshold.
    \item After \texttt{doublings} doublings, the \emph{next} crossing
      \textbf{resets} every bin in that kernel back to \texttt{count}, and
      the cycle repeats.
  \end{itemize}
  \item Example, \texttt{count:1, doublings:6}: report on the 1st, 2nd, 4th,
    8th, 16th, 32nd, 64th occurrence, then reset.
  \item Used throughout Case Studies~1 and~3 to keep hot-kernel logs legible
    while still sampling every order of magnitude of occurrence.
\end{itemize}
\end{frame}

\begin{frame}{Dogfooding NixNan: Two Concurrency Bugs Found \& Fixed}
\label{sl:dogfood}
While porting \texttt{doublings} for this tutorial, two real bugs were found
\textbf{in NixNan itself}:
\begin{enumerate}
  \item \textbf{Doubling-transition race}: \texttt{threshold} and
    \texttt{times\_doubled} are shared, mutable state read-then-written by
    every GPU thread that crosses them --- a classic check-then-act race.
    Evidence: threshold values observed jumping \emph{backward} ($2\to1$
    without ever reaching 4). \textbf{Fix}: gate both transitions with
    \texttt{atomicCAS} so exactly one thread performs each transition.
  \item \textbf{Torn stdout}: the exception-report thread and the
    histogram-report thread both write multi-part \texttt{<<} chains to the
    same stream with no synchronization --- lines from two different reports
    interleave mid-message. \textbf{Fix}: assemble each report as one string,
    write it under a single mutex-protected \texttt{nnout\_line()} call.
\end{enumerate}
\textit{Lesson: exercising a tool's own adaptive/concurrent features hard
enough surfaces real bugs in the tool, not just the target program.}
\end{frame}

\begin{frame}{\texttt{autonixnan}: A 3-Phase Automated Wrapper}
\begin{itemize}
  \item A Python wrapper around \texttt{nixnan.so} --- launches the target
    program multiple times with different env-var combinations, stitches logs.
  \item \textbf{Phase 1 --- baseline}: \texttt{LOG\_KERNELS} mode; catalog
    every kernel and its wall-clock time, \emph{zero} instrumentation overhead.
  \item \textbf{Phase 2 --- extreme-value scan}: \texttt{HISTOGRAM=1},
    exceptions off; find which kernels touch extreme exponents.
  \item \textbf{Phase 3 --- deep dive}: full exception instrumentation,
    optionally whitelisted to the kernels Phase 2 flagged.
  \item Convenient for large, unfamiliar binaries (cuBLAS/cuSOLVER-heavy) ---
    but see the next slide for its main failure mode.
\end{itemize}
\end{frame}

\begin{frame}{A Real Blind Spot: Hard Overflow Jumps}
\label{sl:blindspot}
\begin{itemize}
  \item Case study: a mixed-precision Cholesky solver forced into overflow
    (\texttt{--noscale --bscale 1e5}) --- confirmed by the \emph{program's
    own} internal check (\texttt{nonfinite=256}) on the very first iteration.
  \item \texttt{autonixnan}'s Phase 2 scan \textbf{never flagged the
    responsible kernel} --- it flagged an unrelated FP64 QR kernel used only
    for test-matrix generation. Phase 3's whitelisted deep-dive then reported
    \textbf{zero FP16 exceptions}.
  \item Why: the histogram mechanism's bins exclude zero/infinity by design
    (Slide~\ref{sl:histogram}) --- a one-step jump to \texttt{Inf} leaves no
    gradually-climbing exponent trail for the scan to catch.
  \item \textbf{Fix}: run \texttt{nixnan.so} \textbf{directly}, unwhitelisted,
    no \texttt{autonixnan} --- found the real exceptions immediately.
  \item \textit{Guidance: prefer direct \texttt{nixnan.so} whenever the target
    is small enough to instrument in full.}
\end{itemize}
\end{frame}

% ---------------------------------------------------------------
\section{Real-World Impact}

\begin{frame}{Real-World Impact: the SRU NaN Bug}
\begin{itemize}
  \item \textbf{SRU} (Simple Recurrent Unit): a CUDA implementation of a
    recurrent unit used from PyTorch.
  \item Open issue: \url{github.com/asappresearch/sru/issues/193} --- users
    reported \textbf{NaN outputs} from the library's own example code.
  \item Case study performed by Dr.\ Xinyi Li (GPU-FPX's original author),
    later adapted to run under NixNan.
  \item This is not a synthetic teaching example --- it is a \textbf{real,
    user-filed bug} in a published library.
\end{itemize}
\end{frame}

\begin{frame}[fragile]{SRU Case Study: Root Cause and Fix}
\begin{itemize}
  \item Running the reported-buggy example under NixNan: \textbf{NaN
    exceptions written into memory} appear almost immediately --- far
    upstream of where the user first observed corrupted output.
  \item Following the exception back to source: an \textbf{uninitialized
    tensor}.
\begin{lstlisting}[language=Python]
torch.randn(20, 32, 128).cuda()   # was uninitialized memory before the fix
\end{lstlisting}
  \item Initializing it (instead of leaving it as garbage device memory)
    eliminates the NaNs. Before/after NixNan logs confirm: exceptions
    written into memory before the fix, none after.
  \item \textit{Exactly the pattern this tutorial is about: an exception
    generated far from where a user would ever look, made visible by
    instruction-level instrumentation.}
\end{itemize}
\end{frame}

% ---------------------------------------------------------------
\section{Case Studies}

\begin{frame}{Four Case Studies}
\begin{center}
\begin{tabular}{lll}
\toprule
Case study & Workload & Focus \\
\midrule
\texttt{rd\_nixnan} & Reaction-diffusion PDE solver & FP64 vs.\ FP32 dynamic range \\
\texttt{lu\_solve\_gpu} & LU factorization, partial pivot & Ill-conditioning $\to$ NaN \\
\texttt{tc\_matmult} & Tensor-Core GEMM (WMMA) & Tensor cores + real overhead \\
\texttt{mpir\_gpu} & Mixed-precision iterative refinement & Precision hierarchy at scale \\
\bottomrule
\end{tabular}
\end{center}
\vspace{1em}
All four run on an NVIDIA RTX 3090 / RTX-class Ampere GPU (sm\_80/sm\_86),
instrumented with \texttt{LD\_PRELOAD=nixnan.so}.
\end{frame}

\begin{frame}{Case Study 1: \texttt{rd\_nixnan} --- Reaction-Diffusion}
\[
\frac{\partial u}{\partial t} = D\frac{\partial^2 u}{\partial x^2} + \lambda u,
\qquad \lambda = 25,\ T=4
\]
\begin{itemize}
  \item Growth factor by $t{=}4$: $e^{100}\approx2.7\times10^{43}$ --- inside
    FP64's range, well past FP32's ($\approx3.4\times10^{38}$).
\end{itemize}
\begin{center}
\begin{tabular}{lrr}
\toprule
 & FP64 & FP32 \\
\midrule
Max $|u|$ at $t=3.52$ & $1.11\times10^{38}$ & $1.10\times10^{38}$ \\
Max $|u|$ at $t=3.60$ & $8.09\times10^{38}$ & Overflows to Inf \\
Max $|u|$ at $t=4.00$ & $1.70\times10^{43}$ & NaN (propagated) \\
NixNan exceptions & 0 & 36 NaN / 16 Inf / 40 $-$Inf \\
\bottomrule
\end{tabular}
\end{center}
FP64 completes cleanly; FP32 overflows between $t=3.52$ and $t=3.60$, then NaN
propagates for the rest of the run --- the CPU tutorial's lesson, reproduced on GPU.
\end{frame}

\begin{frame}{\texttt{rd\_nixnan}: A Silent-Failure Wrinkle}
\begin{itemize}
  \item The reduction that reports the running max mirrors
    \texttt{std::max\_element}'s NaN hazard on purpose: \texttt{v > best}
    (not \texttt{fmax}).
  \item Once the array goes non-finite, \emph{every} comparison against
    \texttt{best} fails --- \texttt{best} never leaves its seed value,
    \texttt{0.0} (a boundary point, pinned by the Dirichlet condition).
  \item Result: the program prints a \textbf{plausible, finite \texttt{0.0}}
    from $t{=}3.60$ onward --- not \texttt{nan}, not \texttt{inf}.
  \item This \textbf{silently defeats the program's own}
    \texttt{std::isfinite()} \textbf{self-check} --- it never fires, because
    the reported value genuinely is finite.
  \item NixNan is unaffected: it instruments the FADD/FFMA \emph{inside} the
    step kernel, before the corrupted value ever reaches the reduction. It
    reports the real NaN/Inf (and, with \texttt{INSTR\_MEM}, confirms
    36{,}928 NaN values actually written to global memory) regardless.
\end{itemize}
\end{frame}

\begin{frame}{Case Study 2: \texttt{lu\_solve\_gpu} --- LU Factorization}
$PA = LU$, partial pivoting, forward/backward substitution, $20\times20$
system, ``snap near-epsilon values to zero'' elimination rule.
\begin{center}
\begin{tabular}{lrr}
\toprule
 & \texttt{matrix.csv} (well-cond.) & \texttt{bad\_matrix.csv} (ill-cond.) \\
\midrule
Residual $\lVert Ax-b\rVert_2$ & $9.68\times10^{-16}$ & \texttt{nan} \\
FP64 NaN (repeats) & 0 & 316 (2171) \\
FP64 Infinity & 0 & 4 \\
FP64 Division-by-zero & 0 & 4 \\
\bottomrule
\end{tabular}
\end{center}
\begin{itemize}
  \item Mechanism: a near-zero pivot is snapped to exactly \texttt{0.0} by
    the epsilon-clamp rule; the subsequent division by that zero pivot
    produces $\pm\infty$, and NaN propagates through the rest of the solve.
\end{itemize}
\end{frame}

\begin{frame}{\texttt{lu\_solve\_gpu}: A Binary-Instrumentation False Positive}
\begin{itemize}
  \item The ill-conditioned run also reports \textbf{6 FP32 NaN detections}
    --- despite the solver computing exclusively in \texttt{double}.
  \item Root cause, found via the flagged SASS: \texttt{FFMA R8, RZ, R5, R3}
    --- \texttt{ptxas}'s ``multiply-by-the-zero-register'' idiom for a
    \textbf{register move} ($0{\times}R5+R3=R3$), used to shuffle the two
    32-bit halves of a 64-bit double across the register file.
  \item NixNan correctly classifies \texttt{FFMA} as an \texttt{f32}
    instruction and correctly notices \texttt{R5}'s bit pattern decodes as
    NaN --- but that 32-bit word is \emph{half of a double}, never
    interpreted as an \texttt{f32} value by the program.
  \item \textbf{Takeaway}: a nonzero count in a format the source never uses
    does not necessarily mean that format's arithmetic ran --- a known class
    of false positive for any SASS-level FP exception detector.
\end{itemize}
\end{frame}

\begin{frame}{Case Study 3: \texttt{tc\_matmult} --- Tensor Cores Under NixNan}
\begin{itemize}
  \item WMMA-based FP16/BF16 GEMM (row-major $\times$ row-major), plus
    hand-written SIMT FP32/FP16 comparison kernels.
  \item \textbf{Confirmed empirically}: NixNan's instruction table matches
    tensor-core \texttt{HMMA.16816.F32} directly --- tensor-core MMA
    instructions are \emph{not} invisible to this instrumentation, despite
    operating on packed fragment registers rather than scalar values.
\end{itemize}
\begin{center}
\begin{tabular}{lrr}
\toprule
Kernel & Latency & vs.\ SIMT (same precision) \\
\midrule
FP16, tensor core & $\approx0.178$~ms & \\
FP16, SIMT & $\approx0.865$~ms & $\approx\!4.9\times$ faster (TC) \\
\midrule
BF16, tensor core & $\approx0.177$~ms & \\
FP32, SIMT & $\approx0.845$~ms & $\approx\!4.8\times$ faster (TC) \\
\bottomrule
\end{tabular}
\end{center}
(All $1024^3$, uninstrumented baselines.)
\end{frame}

\begin{frame}{\texttt{tc\_matmult}: Measuring NixNan's Real Overhead}
Four instrumentation levels, one FP16 tensor-core kernel ($128^3$, row-major
$\times$ row-major):
\begin{center}
\scriptsize
\begin{tabular}{lrr}
\toprule
Instrumentation level & Latency & Slowdown \\
\midrule
No NixNan (baseline) & $0.008$~ms & $1\times$ \\
Sampled + histogram (misleading, see next slide) & $0.015$~ms & $\approx\!1.9\times$ \\
Unsampled, default exception checking only & $0.706$~ms & $\approx\!88\times$ \\
Unsampled, exceptions + \texttt{count:1} histogram & $257.009$~ms & $\approx\!32{,}126\times$ \\
\bottomrule
\end{tabular}
\end{center}
\begin{itemize}
  \item Finest-granularity histogramming (\texttt{count:1}) is
    \textbf{drastically} more expensive than exception checking alone ---
    every single occurrence in a watched bin sends a report.
  \item \textit{Instrumentation cost is a real design parameter, not an
    afterthought --- choose granularity deliberately.}
\end{itemize}
\end{frame}

\begin{frame}{\texttt{tc\_matmult}: \texttt{SAMPLING} and a Downsizing Lesson}
\begin{itemize}
  \item The ``sampled'' row above (\texttt{SAMPLING=128}) looks almost free
    ($1.9\times$) --- but this is misleading: this kernel launches only
    $\sim\!88$ times total, so \texttt{SAMPLING=128} instruments only launch
    \#0, never one of the 80 \emph{timed} repeats.
  \item \textbf{Original attempt} at $1024^3$ with \texttt{count:1}: the
    kernel executes $\sim\!262{,}000$ HMMA instructions per launch --- the
    resulting report volume through NixNan's channel meant the run never
    finished in 10+ minutes.
  \item \textbf{Fix}: rerun at $128^3$ ($1/512$ the work) --- tractable in
    under a minute, same code, same watched exponent ranges (re-verified
    empirically for the smaller sample count).
  \item At $128^3$ the tensor-core-vs-SIMT speedup itself collapses to
    $\approx\!1.4\times$ (fixed launch overhead dominates at this size) ---
    a reminder that microbenchmark \emph{size} shapes what you can conclude.
\end{itemize}
\end{frame}

\begin{frame}{Case Study 4: \texttt{mpir\_gpu} --- Mixed-Precision Iterative Refinement}
\begin{itemize}
  \item Carson--Higham iterative refinement Cholesky solve, $n{=}4096$,
    six formats: F64, F32, TF32, F16, BF16-truncated (B16T), BF16-stochastic (B16S).
  \item Ill-conditioned regime ($\kappa$ up to $10^5$): \textbf{all} formats
    break down (\texttt{niter=-1}) --- expected, by design of the test matrices.
  \item Well-conditioned regime ($n{=}256$): convergence-iteration count
    tracks precision directly.
\end{itemize}
\begin{center}
\scriptsize
\begin{tabular}{lcccccc}
\toprule
$\kappa$ & F64 & F32 & TF32 & F16 & B16T & B16S \\
\midrule
$10$ & 0 & 0 & 0 & 2 & 2 & 4 \\
$10^3$ & 1 & 1 & 1 & 3 & 10 & 12 \\
\bottomrule
\end{tabular}
\end{center}
Iterations to reach relative error $<10^{-6}$ --- lower precision needs more refinement steps, especially as $\kappa$ grows.
\end{frame}

\begin{frame}{\texttt{mpir\_gpu}: Subnormal Storms and Extreme Overhead}
Overflow stress test ($\kappa=10^3$, amplified residual $b_{\text{scaled}}=10^5b$),
NixNan-instrumented:
\begin{center}
\scriptsize
\begin{tabular}{lrrr}
\toprule
Format & Baseline & Instrumented & Slowdown \\
\midrule
F16 & 90~ms & 76{,}238~ms & $847\times$ \\
B16T & 155~ms & 79{,}530~ms & $513\times$ \\
F32 & 1.5~ms & 82{,}727~ms & $55{,}152\times$ \\
\bottomrule
\end{tabular}
\end{center}
\begin{itemize}
  \item F16/B16T overhead traces to $\sim$128 \textbf{subnormal (underflow)}
    exceptions per GEMM inside HMMA tensor-core kernels --- F16's 5-bit
    exponent ($\min\approx6{\times}10^{-5}$) is simply too narrow once
    intermediate values are scaled.
  \item F32's overhead is driven by NaN/Inf in QR preprocessing (Householder
    reflections), \emph{not} tensor cores --- a different failure mode
    entirely, same tool catches both.
\end{itemize}
\end{frame}

\begin{frame}{\texttt{mpir\_gpu}: Exceptions Are Not the Same as Failure}
\begin{itemize}
  \item At $\kappa=10^3$, \textbf{every} format shows NixNan-detected
    exceptions (subnormal, div0, NaN, infinity) in QR/factorization ---
    \textbf{yet all still converge}, given enough refinement iterations.
  \item Sweeping $\kappa$ upward: B16T/B16S \textbf{fail} (breakdown,
    \texttt{niter=-1}) at $\kappa\geq2000$; F32/F64 keep converging through
    $\kappa=5000$ (limit of the tested range).
  \item \textbf{Key insight}: the presence of exceptions is not the
    convergence-killer. It is the format's \textbf{precision capacity} ---
    whether iterative refinement has enough headroom to correct the damage
    ---that determines success or failure.
  \item \textit{NixNan tells you exceptions are happening; it does not by
    itself tell you whether that matters for your algorithm's correctness.
    That judgment still belongs to the practitioner.}
\end{itemize}
\end{frame}

% ---------------------------------------------------------------
\section{Lessons and Guidance}

\begin{frame}{Cross-Cutting Lessons for HPC Practitioners}
\begin{itemize}
  \item \textbf{Overflow-then-NaN-propagation} is the single most common
    failure mode across all four case studies --- and it is exactly the
    pattern instruction-level instrumentation is built to catch early.
  \item \textbf{A program's own diagnostics can be fooled}: NaN-hazard
    reductions (\texttt{std::max\_element}-style comparisons) can mask real
    non-finite data behind a plausible finite summary value
    (\texttt{rd\_nixnan}).
  \item \textbf{Instrumentation overhead varies by orders of magnitude}: from
    $\approx\!2\times$ (sparse sampling) to $\approx\!30{,}000\times$
    (unsampled, finest histogram granularity) to $\approx\!55{,}000\times$
    (exception-dense kernels) --- always measure it for your own workload.
  \item \textbf{Binary-level detection has real false positives}
    (\texttt{lu\_solve\_gpu}'s register-move idiom) \textbf{and real blind
    spots} (the histogram scan's zero/Inf exclusion) --- know both before
    trusting a report uncritically.
\end{itemize}
\end{frame}

\begin{frame}[fragile]{Practical Guidance}
\begin{itemize}
  \item \textbf{Start here for a small/medium target}: run
    \texttt{nixnan.so} \textbf{directly}, unwhitelisted, default
    \texttt{INSTRUMENT\_EXCEPTIONS=1} --- no scan-phase blind spot.
\begin{lstlisting}[language=bash]
LD_PRELOAD=./nixnan.so ./target [args...]
\end{lstlisting}
  \item \textbf{For a large/unfamiliar library-heavy binary}: use
    \texttt{autonixnan}'s 3-phase workflow to triage first, but verify its
    Phase-2 kernel selection against a direct run before trusting a
    ``clean'' Phase-3 result.
  \item \textbf{For a hot kernel}: use \texttt{BIN\_SPEC\_FILE} with
    \texttt{doublings} rather than \texttt{count:1} unsampled --- keeps logs
    legible and overhead bounded while still sampling every decade of
    occurrence.
  \item \textbf{Always cross-check} a per-format exception count against
    what the source code actually computes in that format.
\end{itemize}
\end{frame}

\begin{frame}{Conclusion}
\begin{itemize}
  \item NixNan brings instruction-level, no-hardware-trap-required FP
    exception detection to both SIMT cores \emph{and} Tensor Cores on NVIDIA
    GPUs --- including inside closed, precompiled libraries.
  \item Four independent case studies --- a PDE solver, an LU factorization,
    a tensor-core GEMM, and a mixed-precision iterative refinement solver ---
    all reproduce the same central lesson: \textbf{low precision and
    ill-conditioning fail the same way, overflow then NaN propagation, and
    that failure is often invisible until instrumented.}
  \item A real, user-filed library bug (SRU) was found and fixed the same way.
  \item NixNan itself is under active development --- this tutorial's own
    case studies found and fixed two concurrency bugs in the tool.
\end{itemize}
\vspace{1em}
\centering
\url{https://github.com/parfloat/nixnan}
\end{frame}

\end{document}
